\verseline{}{而这一切正在遥远的另一方发生，\footnote{卷首的“而”与“正在发生”把本卷接在第八卷末尾：埃涅阿斯远在伊特鲁里亚盟军一侧受取神盾时，朱诺利用他的缺席，派伊里斯召图尔努斯进攻。卷末伊里斯又奉朱庇特之命传下相反的限制。陶玛斯之女是伊里斯的父称；皮卢姆努斯则是图尔努斯的祖先与拉丁地方神。}}
\verseline{}{萨图尔努斯之女朱诺从天上派下伊里斯，}
\verseline{}{去见勇猛的图尔努斯。他当时恰在圣林里，}
\verseline{}{坐在祖先皮卢姆努斯那奉献于神的谷地。}
\verseline{5}{陶玛斯之女以玫瑰色的唇这样对他说：}
\verseline{}{“图尔努斯，即使你祈求，也没有一位神敢于许诺的事，}
\verseline{}{看啊，岁月轮转，竟自行把它送到你面前。}
\verseline{}{埃涅阿斯抛下城市、同伴和舰队，}
\verseline{}{正前往帕兰忒乌姆，求访厄万德的王廷与居所。}
\verseline{10}{这还不够：他已深入科律图斯的边远城邑，}
\verseline{}{召集吕底亚人的队伍，武装聚来的农人。}
\verseline{}{为何迟疑？现在该唤来战马，现在该唤来战车。}
\verseline{}{再不要耽搁，夺下那陷入惊乱的营地。”}
\verseline{}{她说完，便以双翼升入天空，}
\verseline{15}{从云层下逸去，划出一道巨大的长虹。}
\verseline{}{青年认出了她，便向星辰举起双掌，}
\verseline{}{以呼声追随那远去的身影：}
\verseline{}{“伊里斯，天空的荣光，是谁让云霓载着你，}
\verseline{}{把你送到大地上来？这般明丽的天象}
\verseline{20}{为何骤然显现？我看见天空从中央裂开，}
\verseline{}{星辰在天穹游移。无论是谁召我执起兵器，}
\verseline{}{我都追随这宏大的征兆。”说罢，他走向河边，}
\verseline{}{从河流表面捧取清波，}
\verseline{}{向众神反复祈祷，使高空满是誓愿。}
\verseline{25}{全军已经行进在开阔的原野上，}
\verseline{}{战马众多，彩衣与黄金满目；}
\verseline{}{墨萨波斯统领前锋，提尔鲁斯家的青年统辖后队，}
\verseline{}{主帅图尔努斯行在军阵中央。\footnote{Conte 2019 正文由第 28 行直接接第 30 行；该版校勘记注明传统第 29 行在主要传本中缺失。}}
\verseline{30}{如恒河在七条支流平息后高涨，}
\verseline{}{默然流动；又如水势丰盈的尼罗河，}
\verseline{}{从田野退潮，重新收束进河床。\footnote{本卷第 30—32 行以恒河与尼罗河作双重明喻：恒河以七条已经平静的支流汇成高涨而静默的水势；尼罗河则从泛滥的田野退回河床。两者一盛一收，都写出大军的幅员、力量与有序推进。}}
\verseline{}{特乌克罗斯人望见一团骤起的黑色尘云}
\verseline{}{正在聚拢，阴暗从原野上升起。}
\verseline{35}{卡伊库斯最先从对面的壁垒上高喊：}
\verseline{}{“同胞们，那裹在黑雾里的是什么队伍？}
\verseline{}{快取铁刃，递上投枪，登上城墙——}
\verseline{}{敌人来了，快！”特乌克罗斯人伴着巨大的喊声，}
\verseline{}{穿过所有大门退入营内，布满壁垒。}
\verseline{40}{因为最善战的埃涅阿斯临行前曾这样命令：}
\verseline{}{倘若中途发生任何变故，}
\verseline{}{他们都不可列阵，也不可把希望寄托于原野，}
\verseline{}{只须守住营地和土垒庇护的城墙。}
\verseline{}{因此，尽管羞耻与怒火催他们交手，}
\verseline{45}{他们仍以门闩挡住营门，执行主帅的命令，}
\verseline{}{全副武装，在中空的塔楼里等待敌人。}
\verseline{}{图尔努斯疾驰在前，早已越过迟缓的大队；}
\verseline{}{他只带二十名精选骑兵，猝然来到营城。}
\verseline{}{一匹身带白斑的色雷斯马驮着他，}
\verseline{50}{金盔上的红色羽冠遮护头顶。}
\verseline{}{“青年们，谁愿与我一同率先攻向敌人？}
\verseline{}{看！”他说着，旋臂把标枪掷入空中，}
\verseline{}{以此揭开战斗，昂然冲进原野。}
\verseline{}{同伴以喊声回应，随他向前，}
\verseline{55}{轰然发出可怖的咆哮；他们惊异特乌克罗斯人何以怯懦，}
\verseline{}{何以不肯在平坦的战场迎战，不肯让士兵}
\verseline{}{持兵器正面相逢，只顾紧守营地。图尔努斯狂怒，}
\verseline{}{策马来回察看城墙，在难以通行之处寻找入口。}
\verseline{}{如同一头狼，窥伺着满是羊群的围栏，}
\verseline{60}{饱经风雨，仍在午夜对着羊栏低吼；}
\verseline{}{羊羔安然依偎母羊，不住啼叫，}
\verseline{}{它却凶狠难制，对隔在栏内的猎物狂怒；\footnote{瓦里乌斯·鲁福斯《论死亡》残篇写猎犬追逐已经失去踪影的鹿，其中有 saevit in absentem（“对不在场的猎物狂怒”）；维吉尔在 9.63 几乎原样写成 saevit in absentis，把它转给隔着营垒扑不到特洛伊人的图尔努斯。}}
\verseline{}{积蓄已久的疯狂饥饿折磨着它，}
\verseline{}{咽喉因渴血而干枯：}
\verseline{65}{图尔努斯望着城墙与营地时，也正是如此，}
\verseline{}{怒火燃起，愤痛在坚硬的骨骼中燃烧。}
\verseline{}{怎样才能探得入口？又该循哪条道路，把特乌克罗斯人}
\verseline{}{从封闭的壁垒中震出，使他们倾泻到平原？}
\verseline{}{舰队贴着营地一侧隐蔽停泊，}
\verseline{70}{周围有土垒和河水阻隔；他扑向那里，}
\verseline{}{向欢呼的同伴索要火种，自己也满怀狂热，}
\verseline{}{一把抓满燃烧的松枝。}
\verseline{}{于是众人全力拥上；图尔努斯就在眼前催逼，}
\verseline{}{所有青年都用熏黑的火把武装起来。}
\verseline{75}{他们拆散营火；松脂火炬喷出烟雾与黯光，}
\verseline{}{武尔坎努斯把混杂的火屑送上星空。}
\verseline{}{缪斯啊，是哪位神从特乌克罗斯人那里}
\verseline{}{驱走如此凶猛的火灾？谁从舰船上逐开这样的烈焰？}
\verseline{}{请告诉我。昔人曾信此事；后世长存的却只是传说。}
\verseline{80}{当初，埃涅阿斯在佛里吉亚的伊达山}
\verseline{}{建造舰队，准备驶向远海，}
\verseline{}{据说众神的母亲、柏瑞钦托斯的女神亲自\footnote{“众神的母亲”“柏瑞钦托斯的女神”都指佛里吉亚大母神库柏勒；柏瑞钦托斯是与其崇拜相关的山地称号。本卷追溯伊达山木料所造舰船何以在新的火攻中获得神化，并与第五卷朱诺诱使特洛伊妇女焚船的情节相呼应。第79行把昔人对奇事的信仰与后世长存的传说相对照；第82行“据说”也使叙述保持距离。}}
\verseline{}{以这些话向伟大的朱庇特请求：“孩子，你已平定奥林匹斯，}
\verseline{}{请把亲爱的母亲向你所求的一事赐给我。}
\verseline{85}{我曾有一片松林，多年来为我所爱，}
\verseline{}{是最高山巅的圣林，祭品都送到那里；}
\verseline{}{林中黑色的枞树与枫木巨干使它幽暗。}
\verseline{}{达耳达诺斯的青年需要舰队时，我欣然}
\verseline{}{把这些树木给了他；如今焦灼的恐惧又使我不安。}
\verseline{90}{请消除我的恐惧，也容一位母亲凭祈祷得到这份力量：}
\verseline{}{愿它们不因航行受创，不被任何风暴击败，}
\verseline{}{愿出生在我们的山上，能为它们带来护佑。”}
\verseline{}{那位转动宇宙星辰的儿子回答她：}
\verseline{}{“母亲，你要把命运召向何处？你为这些船索求什么？}
\verseline{95}{凡人之手造出的船体，难道可以享有不死的权利，}
\verseline{}{而埃涅阿斯竟可安然穿过无常的险境？}
\verseline{}{哪位神曾被允许掌有这样大的权力？}
\verseline{}{不过，将来它们完成使命，抵达奥索尼亚港口后，}
\verseline{}{凡从海浪中幸存、把达耳达诺斯统帅}
\verseline{100}{载到劳伦图姆田野的船只，}
\verseline{}{我都会剥去它们凡俗的形体，命它们成为}
\verseline{}{大海的女神，如涅柔斯之女多托}
\verseline{}{与伽拉忒亚，以胸膛划开起沫的海面。”\footnote{朱庇特拒绝让凡人所造的船从一开始便永生，只承诺在它们完成航程、把埃涅阿斯送达劳伦图姆后，使幸存的船脱去凡俗形体。多托与伽拉忒亚是涅柔斯海女，她们的名字说明新神女所属的海中群体。舰船先扯断缆索、以舰首潜水，继而以同等数目的少女形体浮现，变形次序有意保持清楚。}}
\verseline{}{他说完，便以冥府兄弟的斯提克斯河、}
\verseline{105}{以翻涌沥青、黑涡旋转的河岸为誓，}
\verseline{}{点头应允；这一点头使整个奥林匹斯震颤。}
\verseline{}{所以，应许的日子已经到来，帕耳卡三女神}
\verseline{}{也已完成命定的时刻；图尔努斯的侵害}
\verseline{}{提醒母神从神圣的舰船上驱走火把。}
\verseline{110}{此时，一道异样的光首先闪在众人眼前，}
\verseline{}{一片巨云仿佛从曙光女神那里越过天空，}
\verseline{}{伊达山的女神侍从也随之出现；继而一道可怕的声音}
\verseline{}{从空中落下，传遍特洛伊人和鲁图利人的军阵：}
\verseline{}{“特乌克罗斯人，不必惊惶地保卫我的船，}
\verseline{115}{也不必武装双手；图尔努斯可以先焚尽海洋，}
\verseline{}{却绝不能烧毁这些神圣的松木。去吧，挣脱束缚，}
\verseline{}{去吧，成为海上的女神：母神如此命令。”}
\verseline{}{舰船立即各自扯断连接河岸的缆索，}
\verseline{}{仿佛海豚，船首撞角没入水下，}
\verseline{120}{潜向大海深处。继而——惊人的神迹——\footnote{Conte 2019 正文由第 120 行直接接第 122 行；该版校勘记注明传统第 121 行在主要传本中缺失。下一句的“同样多”直接承接前文整支舰队。}}
\verseline{}{此刻便有同样多的少女形体显现，乘波远去。}
\verseline{}{鲁图利人心中惊愕，墨萨波斯本人也受了惊，}
\verseline{}{战马在他手中乱作一团；河流迟疑不前，}
\verseline{125}{台伯河神沙哑作响，从深水中收回脚步。}
\verseline{}{然而，大胆的图尔努斯并没有失去信心；}
\verseline{}{他反倒用言语振奋士气，还主动斥责敌人：}
\verseline{}{“这些奇兆所针对的是特洛伊人；朱庇特亲自}
\verseline{}{夺走了他们惯常的援助；特洛伊人甚至没等到鲁图利人的投枪\footnote{图尔努斯把舰船神化重新解释为朱庇特切断特洛伊人的海上退路，又宣称自己同样拥有命运；他借这种说法在奇迹之后维持军队士气。随后他以“夺妻”类比海伦战争，并以佛里吉亚身份和女性化措辞羞辱特洛伊人。140行先引“毁灭一次足够”的异议，再以从前犯那一次罪本应足够作答。Conte 2019 校勘记在141行将 modo non 解作 nondum（“尚未”），指他们还未彻底憎恶整个女人族类。}}
\verseline{130}{与火焰。海洋从此对特乌克罗斯人封闭，}
\verseline{}{逃亡的希望也已断绝：世界的一半已从他们手中夺走。}
\verseline{}{大地却掌握在我们手里；意大利各族}
\verseline{}{有成千上万人执起兵器。佛里吉亚人夸示的任何}
\verseline{}{神谕和命定答复，都不能使我畏惧。}
\verseline{135}{命运和维纳斯已经得到足够的恩准：特洛伊人}
\verseline{}{确已踏上丰饶的奥索尼亚土地。我也有我的命运：}
\verseline{}{以铁刃铲除这个因夺人妻子而有罪的民族。}
\verseline{}{这种痛苦不只触动阿特柔斯的儿子们，}
\verseline{}{也不只允许迈锡尼拿起兵器。}
\verseline{140}{‘可毁灭过一次，就该足够。’他们从前犯那一次罪，}
\verseline{}{就该足够；他们却还没彻底恨透整个}
\verseline{}{女人族类。如今，他们倚仗中间这道壁垒，}
\verseline{}{倚仗拖延脚步的壕沟——不过是一线隔开死亡的东西——}
\verseline{}{便有了胆量。他们难道没见过特洛伊的城墙——}
\verseline{145}{由涅普顿亲手筑造——也曾沉入烈火吗？}
\verseline{}{然而，你们当中精选的勇士，谁准备以铁刃劈开壁垒，}
\verseline{}{与我一同闯进那惊惶的营地？}
\verseline{}{我攻打特乌克罗斯人，不需要武尔坎努斯铸成的兵器，}
\verseline{}{也不需要一千艘船；尽可让所有伊特鲁里亚人}
\verseline{150}{立刻与他们结盟。特洛伊人不必害怕黑暗和怯懦的偷袭，}
\verseline{}{不必怕我们杀遍城堡各处的守卫、盗走帕拉斯神像，\footnote{Conte 2019 将第 151 行以方括号排印；该版校勘记指出该行由第二卷第 166 行改写而来，并记录其删疑传统。帕拉斯神像即帕拉狄乌姆，是特洛伊城内受弥涅耳瓦庇护的圣像；希腊人夜间盗取它，是特洛伊陷落传说的一环。图尔努斯把盗像、木马与夜袭归入达那俄斯人的“偷窃”，用以衬托自己宣称的白昼正攻。}}
\verseline{}{也不必怕我们藏进木马幽暗的腹中：}
\verseline{}{我已决意在白昼里公开以火围城。}
\verseline{}{我要让他们知道，他们面对的绝非达那俄斯人}
\verseline{155}{和佩拉斯吉青年——赫克托耳把那些人拖到了第十年。}
\verseline{}{如今，白日最好的时辰已经过去，}
\verseline{}{勇士们，便在余下时光为战果欢喜，养足身体，}
\verseline{}{同时等候那正在准备的战斗。”}
\verseline{}{与此同时，墨萨波斯受命以哨兵围住营门，}
\verseline{160}{又以火光环绕城墙。}
\verseline{}{十四名鲁图利人被选出，各守一段城垣；}
\verseline{}{每人又有一百名青年跟随，}
\verseline{}{他们都是紫红羽冠、金光闪烁的青年。}
\verseline{}{众人往来分散，轮换守夜，}
\verseline{165}{散卧草地，纵情饮酒，饮空青铜调酒缸。}
\verseline{}{营火通明；守夜人以嬉游}
\verseline{}{度过整个不眠之夜。}
\verseline{}{特洛伊人从壁垒上望见这一切，}
\verseline{}{全副武装守住高处；他们也因恐惧而警觉，}
\verseline{170}{检查营门，又架起通道，接通各处胸墙，}
\verseline{}{运送投枪。门涅斯透斯和勇烈的塞雷斯图斯临阵督促；}
\verseline{}{父亲埃涅阿斯早已任命他们：危局一旦召来，}
\verseline{}{便统领青年，主持营中事务。}
\verseline{}{全军抽签分担危险，在城墙各处警戒，}
\verseline{175}{轮换值守每个人受命保卫的地段。}
\verseline{}{守门的是尼苏斯，许耳塔库斯之子，精于兵器；}
\verseline{}{善猎的伊达把他送作埃涅阿斯的同伴，}
\verseline{}{他投枪迅疾，射出轻箭也同样迅疾。}
\verseline{}{同伴欧律阿卢斯就在他身旁；在埃涅阿斯族人中，}
\verseline{180}{没有哪个特洛伊青年披上兵甲时比他更美，}
\verseline{}{少年还未生出须发，面容只刻着青春最初的痕迹。}
\verseline{}{两人共有一份爱，也一同奔向战斗；}
\verseline{}{这一次，他们也共同值守同一座营门。\footnote{尼苏斯与欧律阿卢斯已在第五卷祭赛中以深厚的 amor 相连，本卷再让他们共同守门、共同出营。两人的关系同时包含爱、忠诚、求名和愿为对方赴死的决心。}}
\verseline{}{尼苏斯说：“欧律阿卢斯，是神把这股热望放入人心，}
\verseline{185}{还是每个人那可怖的欲望本身都会成为神？\footnote{尼苏斯没有断言自己的热望来自神：他先问是神把冲动放入人心，还是每个人那可怖的欲望本身成为神。}}
\verseline{}{长久以来，我的心总想闯入战斗，或做一件大事，}
\verseline{}{不肯满足于安静、平和的休息。}
\verseline{}{你看鲁图利人对眼前局势何等放心。}
\verseline{}{只有稀疏的火光闪动；酒与睡眠已卸下他们的戒备，}
\verseline{190}{纷纷倒地；四处一片沉寂。再听我说，}
\verseline{}{我正在权衡什么，心中又升起怎样的主意。}
\verseline{}{所有人，无论百姓还是首领，都要求召回埃涅阿斯，}
\verseline{}{派人出去，把确切的消息带回来。}
\verseline{}{若他们答允我替你求取的奖赏——至于我，}
\verseline{195}{这项功业的声名已经足够——我想我能在那座山丘下}
\verseline{}{找到通往帕兰忒乌姆城墙的道路。”}
\verseline{}{欧律阿卢斯怔住，被强烈的荣名之爱攫住；}
\verseline{}{他立即这样对满怀热望的朋友开口：}
\verseline{}{“尼苏斯，难道你竟想避开我，不让我作伴，}
\verseline{200}{参加如此重大的事业？难道我会让你独自走进这般危险？}
\verseline{}{惯经战事的父亲俄斐尔忒斯并非这样养育我，}
\verseline{}{他在阿尔戈斯人的恐怖与特洛伊的苦难间把我抚养成人；}
\verseline{}{与你同行时，我也从来不曾如此行事，}
\verseline{}{那时追随高贵的埃涅阿斯，直至命运的尽头。}
\verseline{205}{我有一颗蔑视生命的心，也相信以生命换取}
\verseline{}{你所奔赴的荣誉，正是相称的代价。”}
\verseline{}{尼苏斯回答：“我从未对你有过这样的担忧，}
\verseline{}{也不该有——宁可伟大的朱庇特，或任何公正观照此事的神}
\verseline{}{不让我凯旋回到你身边，我也绝不会怀疑你的勇气。\footnote{Conte 2019 正文作 nec fas: non ita ...，冒号在 fas 后。该版校勘记另记录在 non 后停顿、从 ita me referat 起誓的读法，并将句意解释为：宁愿朱庇特与自己为敌，也不愿怀疑欧律阿卢斯的勇气。}}
\verseline{210}{然而你也看见，在这样的危局中，变数何其繁多；}
\verseline{}{若有什么意外或神明把我卷向灾祸，}
\verseline{}{我愿你活下去：你的年岁更配活下去。}
\verseline{}{也好有人将我从战斗中夺回，或以赎金换回，}
\verseline{}{交付大地；若命运又照它惯常的方式阻止此事，}
\verseline{215}{也有人能为远方的我献亡者祭品，以空墓致敬。\footnote{尼苏斯设想有人将自己从战场抢回或以赎金换回，以便入土；若尸身无法取回，欧律阿卢斯仍可为远方的他献亡者祭品，以空墓致敬。Conte 2019 在 humo 后断句，将 solita 与 fortuna 相连，指命运又照惯常的方式阻止此事。}}
\verseline{}{我更不愿让你可怜的母亲因我承受如此悲痛；}
\verseline{}{在所有母亲中，少年啊，只有她敢于随你而来，}
\verseline{}{不肯留在伟大的阿刻斯忒斯城中。”}
\verseline{}{欧律阿卢斯却说：“你徒然编织这些空洞的理由，}
\verseline{220}{我的决心不会改变，也绝不从原处退让。}
\verseline{}{让我们快些动身。”说着，他一并唤醒守卫，}
\verseline{}{后继者接过他们的轮值；他离开岗位，}
\verseline{}{亲自伴着尼苏斯前行，两人去寻找王子。}
\verseline{}{大地上其余一切生灵都在睡眠中}
\verseline{225}{卸下忧虑，让心忘记劳苦；}
\verseline{}{特乌克罗斯人的主要首领和精选青年}
\verseline{}{却正为王国最危急的事务举行议事：}
\verseline{}{他们该做什么，谁应立刻给埃涅阿斯送信。}
\verseline{}{众人手持盾牌，倚着长枪站立，}
\verseline{230}{就在营地与原野之间。这时，两名少年一同}
\verseline{}{热切求见，请求获准入内：}
\verseline{}{他们说，此事重大，值得耽搁片刻。尤卢斯首先}
\verseline{}{接纳了忐忑的两人，命尼苏斯开口。}
\verseline{}{许耳塔库斯之子于是说：“埃涅阿斯族人，请以平静的心}
\verseline{235}{听我们陈述；不要因我们的年岁判断}
\verseline{}{我们带来的主意。鲁图利人已被酒与睡眠埋没，}
\verseline{}{沉寂无声；我们亲眼看见一处可供偷越的地方，}
\verseline{}{就在最靠海的营门外，那条岔路敞开着；}
\verseline{}{营火已经断续，黑烟仍升向星辰。}
\verseline{240}{若你们准许我们利用时机，去寻找埃涅阿斯}
\verseline{}{和帕兰忒乌姆城，不久便会看见我们回到此地，}
\verseline{}{携着战利品，已经完成一场巨大的屠杀。}
\verseline{}{道路也不会欺骗我们：}
\verseline{}{我们曾在幽暗山谷下望见城市最外侧，}
\verseline{245}{又在不断打猎时熟悉了整条河流。”}
\verseline{}{此时，年岁深重、心智成熟的阿勒忒斯说：}
\verseline{}{“祖国诸神啊，特洛伊始终立于你们的神威之下，}
\verseline{}{你们终究没有准备把特乌克罗斯人彻底消灭，}
\verseline{}{否则不会把这样的心志、这样坚定的胸怀}
\verseline{250}{赐给我们的青年。”他说着，握住两人的肩膀}
\verseline{}{和右手，泪水流过面容，浸湿双颊。}
\verseline{}{“勇士们，你们配得什么，我又能想出什么}
\verseline{}{奖赏，来报偿这等功业？最美的酬报，}
\verseline{}{首先将由众神和你们自身的品格给予；其余一切，}
\verseline{255}{虔敬的埃涅阿斯很快会偿还；阿斯卡尼乌斯正值青春，}
\verseline{}{也永远不会忘记这样大的功绩。”}
\verseline{}{“不，应当由我——我的救援只在父亲归来——”}
\verseline{}{阿斯卡尼乌斯接着说，“尼苏斯，凭伟大的家神、\footnote{尤卢斯与阿斯卡尼乌斯是同一人。第 257 行开始的话延续到誓言结束；中间“阿斯卡尼乌斯接着说”是叙述插语，仍由同一位说话者继续。}}
\verseline{}{凭阿萨剌科斯的守护神与白发维斯塔的内殿，}
\verseline{260}{我恳求你们；我所拥有的时势与信义，}
\verseline{}{全都托在你们怀中。把我的父亲召回来，}
\verseline{}{让他重回我眼前；他一归来，便再无悲苦。}
\verseline{}{我会给你们两只白银制成、浮雕凸起的酒杯，}
\verseline{}{那是父亲攻陷阿里斯巴时夺得的战利品；}
\verseline{265}{另有两只三足鼎、两大塔兰同黄金，}
\verseline{}{以及西顿的狄多赠送的古老调酒缸。}
\verseline{}{倘若我父亲以胜利者夺取意大利，掌握王权，}
\verseline{}{并有机会抽签分配战利品——}
\verseline{}{你见过图尔努斯所乘的战马，也见过他遍体黄金}
\verseline{270}{行军时的兵甲；那匹马、盾牌和赤红羽冠}
\verseline{}{我会先从抽签中取出：尼苏斯，它们从此就是你的奖赏。}
\verseline{}{此外，父亲将给你们十二名精挑的女俘、}
\verseline{}{男俘和每个人各自的武装，}
\verseline{}{再加上拉提努斯王亲自占有的那片土地。\footnote{银杯、三足鼎、黄金、调酒缸、图尔努斯的战马与兵甲、土地，以及被俘的妇女和其他俘虏，都列在胜利后的奖赏中。俘虏与器物、牲畜和土地并列，说明活人也被当作可占有的财产。}}
\verseline{275}{至于你，可敬的少年，我们年岁相去不远；}
\verseline{}{我已全心接纳你，}
\verseline{}{拥你作伴，与你共赴一切际遇。}
\verseline{}{无论有何作为，我都不会撇下你去求取荣耀；}
\verseline{}{无论我经营和平还是战争，我都会把最高的信任}
\verseline{280}{交给你，无论在行动还是言语上。”欧律阿卢斯回答他：}
\verseline{}{“任何时日都不会证明我不配这般勇敢的事业，}
\verseline{}{无论命运降下顺境还是逆境。}
\verseline{}{但在一切礼物之上，我只求你一件事：}
\verseline{}{我的母亲出自普里阿摩斯的古老族系；我离开时，}
\verseline{285}{可怜的她既没有被伊利乌姆的土地留下，}
\verseline{}{也没有留在阿刻斯忒斯王的城中。}
\verseline{}{如今，我要离开她，让她不知道前方有怎样的危险，}
\verseline{}{也不向她告别；黑夜与你的右手都可以作证：}
\verseline{}{我无法承受母亲的泪。求你安慰}
\verseline{290}{这无依的女人，援助被留下的她。}
\verseline{}{请容我带着对你的这份希望；我将更勇敢地}
\verseline{}{走进一切危局。”达耳达诺斯人心被刺中，}
\verseline{}{流下眼泪；美丽的尤卢斯尤其如此，}
\verseline{}{这幅孝亲的图景触动了他的心。}
\verseline{295}{随后，他这样开口：}
\verseline{}{“你可以相信：一切都将配得上你宏大的开端。}
\verseline{}{她将是我的母亲，只缺少克瑞乌萨这个名字；}
\verseline{}{她生下这样的儿子，绝不会只得到微薄的感激。}
\verseline{}{无论什么际遇跟随这次行动，}
\verseline{300}{我都以这颗头起誓——父亲过去也常以它起誓：}
\verseline{}{我向平安归来、事业顺利的你所许诺的一切，}
\verseline{}{也会同样留给你的母亲和家族。”\footnote{欧律阿卢斯的母亲出自普里阿摩斯族系，特洛伊覆亡后没有留在故土，也没有留在阿刻斯忒斯的西西里城市，随儿子来到拉提乌姆。阿斯卡尼乌斯答允视她为母，只缺“克瑞乌萨”之名，并保证无论此行结局如何，许给儿子的奖赏仍留给她和家族。克瑞乌萨是阿斯卡尼乌斯在特洛伊失去的母亲。}}
\verseline{}{他说着仍在流泪，同时从肩上解下一柄镀金的剑，送给欧律阿卢斯；}
\verseline{}{克诺索斯人吕卡翁曾以奇巧手艺把它造就，}
\verseline{305}{又为它配上贴合的象牙剑鞘。}
\verseline{}{门涅斯透斯把一张带鬃毛的狮皮}
\verseline{}{送给尼苏斯；忠实的阿勒忒斯则同他交换头盔。}
\verseline{}{两人全副武装，立即启程；他们走时，}
\verseline{}{所有首领，连同青年和老人，护送他们}
\verseline{310}{来到营门，以誓愿相随。美丽的尤卢斯也在其中；}
\verseline{}{他的志气与男子的忧思远远超越年龄，}
\verseline{}{托付许多话，要他们带给父亲；可那些话}
\verseline{}{全被风撕碎，徒然交给云层。\footnote{阿斯卡尼乌斯尚未成年，却已经怀着成人统帅的忧虑，把许多话托给两位青年带给父亲。紧接着，风把这些话撕散、交给云层，明说口信不会抵达；这一细节既写出阿斯卡尼乌斯开始承担统帅责任，也再次强调埃涅阿斯仍不在营中。}}
\verseline{}{他们出营，越过壕沟，穿过夜的阴影，}
\verseline{315}{走向敌营；可在自己死去以前，他们将先给许多人}
\verseline{}{带来毁灭。人们因酒与睡眠横卧草间，}
\verseline{}{岸边的战车向上支着车辕；}
\verseline{}{士兵睡在缰绳与车轮之间，兵器与酒散落一处。}
\verseline{}{许耳塔库斯之子首先开口：}
\verseline{320}{“欧律阿卢斯，我们必须用右手冒险；时势正在召唤。}
\verseline{}{道路就在这里。你守望远处，仔细察看，}
\verseline{}{别让任何人从背后向我们举起手来；}
\verseline{}{我会使这一带荒芜，为你铺出一条宽路。”}
\verseline{}{他说完压低声音，同时持剑直扑拉姆涅斯；}
\verseline{325}{这傲慢的人恰好高高躺在一堆毯子上，}
\verseline{}{鼾声从整个胸膛沉重地涌出。}
\verseline{}{他是国王，也是图尔努斯王最钟爱的占卜者，}
\verseline{}{却没能凭占卜驱走这场灾难。}
\verseline{}{尼苏斯又杀死近旁三个仆人，他们胡乱睡在兵器间；}
\verseline{330}{又找到雷穆斯的执盾者与御者；他们卧在战马下，}
\verseline{}{垂落着颈项，尼苏斯用铁刃一一割断；}
\verseline{}{接着，他砍下主人本人的头颅，把无首的躯干留下，}
\verseline{}{让它抽噎着喷血；土地与床铺}
\verseline{}{全被温热的黑血浸透。他又杀了拉弥鲁斯、拉穆斯，}
\verseline{335}{还有年轻的塞拉努斯；他容貌出众，}
\verseline{}{那一夜嬉游良久，此刻四肢倒下，}
\verseline{}{深深被酒神制服——倘若他一刻不停，}
\verseline{}{把嬉游延续整夜、直至晨光，反倒会是幸运的。}
\verseline{}{如一头饥饿的狮子闯入拥挤的羊栏，}
\verseline{340}{疯狂的饥饿驱使它，啃噬、拖走}
\verseline{}{柔弱而因恐惧失声的羊群，血口低吼：}
\verseline{}{欧律阿卢斯的杀戮也不逊色；他同样燃烧起来，}
\verseline{}{狂怒地穿行，在人群中央杀死大批无名之人；}
\verseline{}{又逼近法杜斯、赫耳柏苏斯、罗厄图斯和阿巴里斯，}
\verseline{345}{他们全不知情；唯有罗厄图斯醒着，看见一切，}
\verseline{}{却因恐惧躲在一只巨大的调酒缸后。}
\verseline{}{他起身时，欧律阿卢斯迎面把整柄剑}
\verseline{}{近身埋入胸膛，又把沾满死亡气息的剑拔出。}
\verseline{}{那人吐出紫红的生命，临死时把混着血的酒}
\verseline{350}{一并呕出；欧律阿卢斯却仍在偷袭的狂热中进逼。}
\verseline{}{他已走向墨萨波斯的同伴；在那里，最后一簇火光}
\verseline{}{正在熄灭，按序系好的战马低头啮草。}
\verseline{}{尼苏斯这时简短地开口；}
\verseline{}{他已觉出，欧律阿卢斯被过度的杀戮与欲望卷走：}
\verseline{355}{“收手吧，敌对的晨光正在逼近。}
\verseline{}{报复已经足够，穿过敌阵的道路也已打开。”\footnote{二人趁敌军醉眠夜入营地，取用了《伊利亚特》第十卷奥德修斯与狄俄墨得斯夜袭特洛伊盟军的故事；第一卷迦太基神庙壁画中的瑞索斯之死，正属于那次荷马夜袭。维吉尔却让自己的夜袭者在返途中被发现并双双死去。}}
\verseline{}{他们留下将士们的许多纯银精制器物，}
\verseline{}{连同兵器、调酒缸和美丽的毯子。}
\verseline{}{欧律阿卢斯取走拉姆涅斯的胸饰和金钉佩带；}
\verseline{360}{它们属于提布尔人雷穆卢斯；巨富凯狄库斯昔日}
\verseline{}{远隔两地同他缔结宾主之谊时，把这些作为礼物送来；}
\verseline{}{雷穆卢斯临死，又把它们留给自己的孙子；}
\verseline{}{孙子死后，鲁图利人才在战争与交锋中夺得这些物件。\footnote{拉姆涅斯的胸饰与金钉佩带原由凯狄库斯远赠提布尔人雷穆卢斯，雷穆卢斯留给孙子，孙子战死后才落入鲁图利人手中，如今又被欧律阿卢斯夺走。胸饰即 phalerae，可指佩在胸甲上的金属饰牌；塞尔维乌斯（9.359）在此另取马具饰物。装备由宾主赠礼到家族传承，再因战死转手；墨萨波斯头盔的反光随后暴露了两人的行踪。}}
\verseline{}{欧律阿卢斯把它们抢来，徒然佩在强壮的肩上，}
\verseline{365}{又戴上墨萨波斯那顶合身而饰有羽冠的头盔。}
\verseline{}{两人离开营地，向安全之处走去。}
\verseline{}{与此同时，从拉丁城先行出发的骑兵}
\verseline{}{正在赶路；其余军团列阵停在原野，}
\verseline{}{他们正把国王的答复带给图尔努斯，}
\verseline{370}{共三百人，全都持盾，由沃尔肯斯统领。\footnote{沃尔肯斯所率三百骑兵由拉丁城先行出发，是外来报信队伍；他们不同于营内醉睡的守军。欧律阿卢斯的头盔在微明夜色中迎光闪耀，骑兵才封锁熟悉的森林岔路。}}
\verseline{}{他们已经靠近营地，正在走向城墙，}
\verseline{}{忽从远处看见两人折向左边的小径；}
\verseline{}{微明夜色的阴影中，欧律阿卢斯的头盔}
\verseline{}{泄露了忘却危险的他，迎着光芒闪耀。}
\verseline{375}{这景象立刻引起警觉。沃尔肯斯从队伍中高喊：}
\verseline{}{“站住，勇士们！你们为何上路？究竟是谁，竟全副武装？}
\verseline{}{要往哪里去？”两人没有回话，}
\verseline{}{只是加快脚步逃向森林，把希望寄托给黑夜。}
\verseline{}{骑兵熟知岔道，从两侧}
\verseline{380}{挡住去路，以守卫包围每一个出口。}
\verseline{}{那里有一片宽广的森林，灌木与黑栎使它幽暗，}
\verseline{}{又遍布粗硬的密集荆棘；}
\verseline{}{只有一条隐约的小径，在密林暗径间透出微光。}
\verseline{}{枝叶的黑暗和沉重的战利品阻碍欧律阿卢斯，}
\verseline{385}{恐惧又使他认错道路的方向；}
\verseline{}{尼苏斯已经走远，不知不觉逃出了敌阵，}
\verseline{}{来到后来因阿尔巴而得名的地方；那里日后称作阿尔巴尼，}
\verseline{}{而当时拉提努斯王还在那里拥有高大的畜栏。}
\verseline{}{他终于停下，徒然回望不在身旁的朋友。}
\verseline{390}{“不幸的欧律阿卢斯，我把你丢在了哪里？\footnote{Conte 2019 在 Euryale infelix 之后断句，并注明这一断句带有保留；按正文句读，infelix 连到欧律阿卢斯。另一古老断法在 Euryale 后停顿，可把 infelix 理解为尼苏斯自称“不幸”。}}
\verseline{}{我该沿哪条路去找？”他重新穿过这座欺人的森林，}
\verseline{}{把每一条纠结的路径重走一遍；说着，他循着}
\verseline{}{先前记下的足迹折返，在沉默的灌木间游走。}
\verseline{}{他听见战马，听见追兵的脚步与信号。}
\verseline{395}{没有多久，呼喊声便传到耳边；}
\verseline{}{他看见欧律阿卢斯已被整支队伍压制，}
\verseline{}{地势与黑夜蒙骗了他，骤起的喧嚣使他惊乱，}
\verseline{}{他正被拖走，虽百般挣扎也全然徒劳。}
\verseline{}{他该怎么办？凭何种力量、何种兵器，}
\verseline{400}{才敢夺回少年？还是将自己投进包围，}
\verseline{}{奔入剑阵求死，借伤口赶赴美丽的死亡？}
\verseline{}{他迅速引臂向后，转动手中的枪杆，}
\verseline{}{仰望高悬的月神，以这样的声音祈祷：}
\verseline{}{“女神啊，请你亲自临在，援助我们的困苦；}
\verseline{405}{星辰的荣光，拉托娜之女，林野的守护者，}
\verseline{}{若我父亲许耳塔库斯曾为我把礼物}
\verseline{}{带到你的祭坛，若我自己也从狩猎所得为你增添祭品，}
\verseline{}{把它们悬上圆顶，或钉在神殿的圣檐，}
\verseline{}{请准我扰乱这支队伍，引导我的投枪穿过空气。”\footnote{尼苏斯祈求的月神又被称作“拉托娜之女”和林野守护者，即狄安娜。他以自己和父亲许耳塔库斯献上的狩猎祭物为依据，请女神引导投枪；随后从隐处杀死苏尔摩和塔古斯。}}
\verseline{410}{他说完，以全身力量掷出铁枪；}
\verseline{}{飞枪劈开黑夜的阴影，}
\verseline{}{击中背向他的苏尔摩后背，在那里折断，}
\verseline{}{裂开的木杆却仍穿过胸中要害。}
\verseline{}{那人翻滚，胸口吐出温热的血河，}
\verseline{415}{身体渐冷，腰腹被绵长的抽噎反复震动。}
\verseline{}{众人四面张望。尼苏斯因此更加勇猛，}
\verseline{}{看啊，又在耳际高高平衡另一支投枪。}
\verseline{}{敌人仍在惊惶，枪已嘶鸣着穿过}
\verseline{}{塔古斯两边的太阳穴，钉在脑髓里，渐渐温热。}
\verseline{420}{残暴的沃尔肯斯狂怒，却看不见投枪}
\verseline{}{从何处发出，也找不到狂怒扑击的方向。}
\verseline{}{“那么在此期间，你就用温热的血，}
\verseline{}{替他们两人向我偿命。”他说着拔出剑，}
\verseline{}{扑向欧律阿卢斯。尼苏斯这才惊骇失神，}
\verseline{425}{放声大喊；他再也不能藏身黑暗，}
\verseline{}{也无法承受如此巨大的痛苦：}
\verseline{}{“向我来，向我来！我在这里，是我做的！把铁刃}
\verseline{}{转向我，鲁图利人！一切计谋都是我的；这个少年}
\verseline{}{什么也没敢做，什么也没能做；凭这天空和知情的星辰作证，}
\verseline{430}{只是太爱我这个不幸的朋友。”}
\verseline{}{他这样呼喊；可是沃尔肯斯已经发力刺出长剑，}
\verseline{}{剑锋穿过肋骨，破开洁白的胸膛。}
\verseline{}{欧律阿卢斯翻滚着坠入死亡，鲜血流过}
\verseline{}{美丽的四肢，颈项垂落，倚在肩上；}
\verseline{435}{如同一朵紫红的花被犁刃削断，}
\verseline{}{垂死萎谢；又如罂粟以疲倦的茎颈}
\verseline{}{低下花首，偶然被雨水压重。}
\verseline{}{尼苏斯却直冲人群，越过所有敌手，}
\verseline{}{只扑向沃尔肯斯，眼中只剩沃尔肯斯一人。}
\verseline{440}{敌人从两面紧紧围拢，近身}
\verseline{}{把他逼退；他仍不停下，旋动}
\verseline{}{闪电般的剑，直到把它迎面埋进}
\verseline{}{高喊的鲁图利人嘴中，自己临死也夺去敌人的生命。}
\verseline{}{随后，他遍体被刺，扑在死去的朋友身上，}
\verseline{445}{终于在那里安息于宁静的死亡。\footnote{尼苏斯为救朋友，主动现身并把全部计谋揽在自己身上；这是一项自我牺牲的谎言。欧律阿卢斯之死先比作被犁割断的紫花，又比作受雨压重、因颈项疲倦而低头的罂粟；柔软、重力与渐萎紧接在穿胸的暴力之后。尼苏斯杀死沃尔肯斯后扑在朋友尸体上，也在那里死去。}}
\verseline{}{幸运的二人啊！只要我的诗歌还有力量，}
\verseline{}{任何时日也不能把你们从世代记忆中夺走，}
\verseline{}{只要埃涅阿斯的家族仍居于卡庇托利乌姆不动的岩石，}
\verseline{}{罗马之父仍握有统治权。\footnote{“幸运的二人啊”是诗人直接呼告；诗歌许诺他们长存于世代记忆。“罗马之父”（pater Romanus）可联系奥古斯都、朱庇特或罗马家长的权威。埃涅阿斯家族、卡庇托利乌姆与这一父亲称谓把私人情感接入罗马未来。紧接这段不朽声名的许诺，诗中又写二人的头颅被割下示众。}}
\verseline{450}{获胜的鲁图利人夺得掠获物与剥下的兵甲，}
\verseline{}{却哭泣着把死去的沃尔肯斯抬回营地。}
\verseline{}{营内发现拉姆涅斯血已流尽，}
\verseline{}{又发现这么多首要人物同时被杀，悲痛并不稍减；}
\verseline{}{塞拉努斯和努玛也在其中。人群一齐聚拢，}
\verseline{455}{冲向尸体和尚未断气的人，也冲向尚有余温、}
\verseline{}{屠杀方歇的土地与满是泡沫血液的水沟。}
\verseline{}{他们彼此认出战利品：墨萨波斯闪亮的头盔，}
\verseline{}{还有历经辛劳才夺回的胸饰。}
\verseline{}{初晓时，奥罗拉正离开提托诺斯藏红色的床榻，}
\verseline{460}{以新生的光洒遍大地；}
\verseline{}{阳光已经倾泻，万物已在光中显露，}
\verseline{}{图尔努斯自己披满兵甲，唤起战士执起武器，}
\verseline{}{又唤起披铜甲的战阵；各统帅都驱策部下}
\verseline{}{投入战斗，用纷纭的传言磨锐怒火。\footnote{Conte 2019 在 aeratasque acies 后断句，把 uiros ... aeratasque acies 同作图尔努斯 suscitat 的宾语；随后 in proelia cogit / quisque suos 另成句。因此披铜甲的战阵仍是图尔努斯“唤起”的对象，quisque suos 才转指各统帅各自的人。}}
\verseline{465}{甚至——目不忍视——他们把两颗头颅}
\verseline{}{钉上高举的枪尖，在震天喊声中前进：}
\verseline{}{那是欧律阿卢斯与尼苏斯。}
\verseline{}{坚忍的埃涅阿斯族人在城墙左侧}
\verseline{}{列阵迎敌，右侧则由河流环绕；}
\verseline{470}{他们守住巨大的壕沟，站在高塔上，}
\verseline{}{满怀哀伤；那些插在枪上的面容使他们动摇，}
\verseline{}{可怜的人们太熟悉它们，黏黑的血水正从上面流下。\footnote{鲁图利人一面为沃尔肯斯、拉姆涅斯等死者哀悼，一面把尼苏斯和欧律阿卢斯的头颅插在枪尖上示众。黎明把夜间暗杀的结果转为双方共同观看的公共景象；黑色腐血与可辨认的面孔又把诗人刚许诺的永恒声名拉回残损肉身。}}
\verseline{}{与此同时，有翼的传闻飞过惊惶的营城，}
\verseline{}{这报信者猛冲而来，落到}
\verseline{475}{欧律阿卢斯母亲耳边。霎时，温热从这可怜女人的骨中退尽，}
\verseline{}{梭杆从手中震落，纺出的线团也散开。}
\verseline{}{不幸的她飞奔而出，发出女人的号哭，}
\verseline{}{撕扯头发，狂乱地冲向城墙和前列军阵；}
\verseline{}{她顾不得战士，顾不得危险，}
\verseline{480}{也顾不得投枪，只以哀号填满天空：}
\verseline{}{“欧律阿卢斯，我看见的竟是这样的你？你——我迟暮}
\verseline{}{之年唯一的慰藉——竟忍心独自留下我，}
\verseline{}{残酷的孩子？你被派往如此巨大的危险，}
\verseline{}{可怜的母亲竟连最后同你说话也未获准许？}
\verseline{485}{唉，你躺在陌生的土地上，成了拉丁人的狗}
\verseline{}{与飞鸟的猎物；我这母亲未能为你送葬，}
\verseline{}{未能为你合上双眼、清洗伤口，}
\verseline{}{也没能用衣袍覆盖你——我日夜赶织那件衣袍，}
\verseline{}{用织机安抚衰年无尽的忧思。}
\verseline{490}{我该到哪里追随你？如今哪片土地藏着你的躯体、}
\verseline{}{扯断的四肢和破碎的尸首？孩子，}
\verseline{}{你只把这些带回给我？我随你渡海越陆，只为了这些？}
\verseline{}{若还有一点敬义，就刺穿我吧；把所有投枪}
\verseline{}{都掷向我，鲁图利人！先用铁刃夺去我的生命。}
\verseline{495}{或者，伟大的众神之父，请你怜悯，}
\verseline{}{以雷霆把这可恨的头打入塔耳塔洛斯，}
\verseline{}{因为除此以外，我无法割断这残酷的生命。”}
\verseline{}{她的哭声震动众人的心，悲哀的呻吟}
\verseline{}{传遍军队；战士们战意折断，力量凝滞。}
\verseline{500}{伊达乌斯和阿克托耳见她还在煽起众人的悲痛，}
\verseline{}{便依伊利俄纽斯和泪流不止的尤卢斯之命，}
\verseline{}{一把抓住她，将她抱起，送回屋内。\footnote{这里的 Fama 生翼飞行，具有报信者的拟人化力度。死讯使母亲手中梭杆坠落、纺线散开；她哀叹自己不能为儿子合眼、洗伤、覆衣、送葬，也不知道四肢和残身落在何处。她的哭声使守军战意低落，伊达乌斯和阿克托耳奉命把她抱回屋内。马克罗比乌斯（6.2.21）又把“未能为你送葬、合上双眼、清洗伤口”与恩尼乌斯失传悲剧《克瑞斯丰忒斯》（Cresphontes）中的母亲哀词并列：后者同样痛诉无缘覆土、覆尸和以泪洗血。}}
\verseline{}{远处，号角的铜腔迸出}
\verseline{}{可怕的声响；喊声紧随而起，天空回鸣。}
\verseline{505}{沃尔斯基人急速推进，共同举起龟甲阵，\footnote{testudo 指士兵以盾牌在头顶和四周相覆形成的龟甲阵，用于接近壕沟和壁垒。特洛伊人先以巨石砸散盾顶，鲁图利人转用投射兵器；墨泽提乌斯负责送来火焰，墨萨波斯则拆除壁障、索取梯子。}}
\verseline{}{准备填平壕沟，拔除壁垒。}
\verseline{}{一部分人寻找入口，以梯子攀登城墙，}
\verseline{}{攻向阵线稀疏、守军环列间透出空隙的地方，}
\verseline{}{那里人群并不密集。特乌克罗斯人倾泻}
\verseline{510}{一切投射兵器，又以坚硬长杆把他们推下；}
\verseline{}{漫长的战争早已教会他们怎样保卫城墙。}
\verseline{}{他们还滚下沉重而来势凶猛的岩石，企图}
\verseline{}{砸破盾牌遮护的军阵；可盾阵中的士兵}
\verseline{}{仍甘愿在密集的龟甲下承受一切打击。}
\verseline{515}{可盾阵终于支撑不住。敌群大队逼近之处，}
\verseline{}{特洛伊人把一块巨石}
\verseline{}{滚下、推落；它使鲁图利人大片倒下，}
\verseline{}{拆散了兵甲相覆的顶盖。大胆的鲁图利人}
\verseline{}{不再愿意在幽暗的近战中相争，}
\verseline{520}{转而竞相以投射兵器把守军赶下壁垒。}
\verseline{}{另一侧，可怖的墨泽提乌斯挥动}
\verseline{}{一根伊特鲁里亚松木，把浓烟滚滚的火焰送来；}
\verseline{}{而墨萨波斯，驯马者、涅普顿之子，}
\verseline{}{正在拆开壁垒，高声召来登城的梯子。}
\verseline{525}{卡利俄佩啊，请你启发我的歌声，}
\verseline{}{告诉我图尔努斯当时以铁刃造成多少毁灭、多少死亡，}
\verseline{}{各人又把哪一名勇士送进奥尔库斯；}
\verseline{}{请与我一同铺开这场大战的辽阔场面，}
\verseline{}{因为女神们，你们记得，也有力量把它讲述。\footnote{第525行呼名卡利俄佩，却用复数 vos、adspirate；塞尔维乌斯（9.525）指出这里合用了单数名号与复数呼求。第529行再称“女神们”。Conte 2019 将第529行置于方括号；该版校勘记记录该行重现第七卷645行，并在若干主要传本中缺失。}}
\verseline{530}{营中有一座塔楼，高得令人仰望，木台高架，}
\verseline{}{位置有利；全体意大利人正竭尽全力，}
\verseline{}{动用最大兵力，争相攻下它、将它推倒；}
\verseline{}{特洛伊人则以岩石防守，守军密集，}
\verseline{}{从空阔的窗洞中掷出兵器。}
\verseline{535}{图尔努斯最先投出一支燃烧的火把，}
\verseline{}{把火焰钉上塔楼侧面；烈焰乘风攫住}
\verseline{}{大片木板，又咬住被烧蚀的门柱。}
\verseline{}{塔内众人惊乱奔逃，却徒然想从灾难中}
\verseline{}{找出退路；当他们聚成一团，退往}
\verseline{540}{尚未受灾的一面，塔楼骤然因重量}
\verseline{}{向前倒塌，整个天空都因巨响而震鸣。}
\verseline{}{半死的人坠向地面，巨塔的重压紧追而下；}
\verseline{}{他们被自己的投枪刺穿，胸膛也被}
\verseline{}{坚硬的木料贯透。只有赫勒诺耳与吕库斯}
\verseline{545}{勉强逃出。两人之中，赫勒诺耳年纪尚轻；}
\verseline{}{奴女利库姆尼娅曾为迈俄尼亚王偷偷生下他，}
\verseline{}{又把他送往特洛伊，佩着本不许他佩带的兵器；}
\verseline{}{他轻装持裸剑，携素白圆盾，尚未赢得荣名。}
\verseline{}{当他看见自己身处图尔努斯成千上万人中央，}
\verseline{550}{两边都是列阵而立的拉丁军队，}
\verseline{}{如同一头野兽，被密集的猎人围成一圈，}
\verseline{}{仍向投枪狂奔，明知死亡，却把身体投向它，}
\verseline{}{跃起越过猎矛；这青年也正是如此，}
\verseline{}{准备赴死，冲进敌军中央，向兵器}
\verseline{555}{最密集之处直奔而去。}
\verseline{}{吕库斯的双脚则远为迅捷；他在兵器与敌人之间}
\verseline{}{逃亡，直奔高墙，竭力伸手}
\verseline{}{去抓高处墙垛，触到同伴伸出的右手。}
\verseline{}{图尔努斯以脚力和投枪同时追逼，}
\verseline{560}{获胜后这样责骂他：“疯子，你竟以为}
\verseline{}{能够逃出我的手掌？”说着，他抓住}
\verseline{}{悬在墙上的吕库斯，把他连同一大片城墙扯下。}
\verseline{}{如同朱庇特的执兵之鸟以弯曲的利爪}
\verseline{}{高飞，提走一只野兔或白身天鹅；}
\verseline{565}{又如玛尔斯之狼从羊栏夺走}
\verseline{}{母羊不住啼叫寻找的羔羊。四面喊声升起；}
\verseline{}{敌人闯入，动手填平壕沟，}
\verseline{}{另一些人把燃烧的火把掷向屋顶。\footnote{图尔努斯投火后，塔内守军聚到尚未燃烧的一侧，塔楼因集中重量倒塌，半死者又被自己的兵器、硬木和塔体穿刺。逃出的赫勒诺耳主动冲进最密集的兵器，吕库斯则试图攀墙返营。前者死在图尔努斯包围中，后者被图尔努斯从墙上抓下；随后分别出现鹰夺猎物和狼夺羔羊的比喻。530行 pontibus 可指塔楼的木平台或楼层，塞尔维乌斯（9.530）释作 tabulata。}}
\verseline{}{伊利俄纽斯举起一块巨岩、山体的庞大碎片，}
\verseline{570}{击倒正走向营门、带着火焰的琉刻提乌斯；}
\verseline{}{利革尔杀死厄玛提翁，阿西拉斯杀死科律奈乌斯，}
\verseline{}{一个善用标枪，一个善射从远处出其不意袭来的箭；}
\verseline{}{凯纽斯杀死俄耳提癸乌斯，图尔努斯又杀死获胜的凯纽斯。}
\verseline{}{图尔努斯还杀了伊堤斯、克洛尼乌斯、狄俄克西普斯、普罗摩卢斯，}
\verseline{575}{萨伽里斯，以及站在最高塔前的伊达斯；}
\verseline{}{卡皮斯则杀死普里韦耳努斯。特弥拉斯的轻枪最初}
\verseline{}{只擦伤他；他疯狂地扔下盾牌，}
\verseline{}{伸手摸向伤口；于是卡皮斯的箭张翼滑来，}
\verseline{}{把那只手钉在左肋上，深入体内，}
\verseline{580}{以致命伤口割断了灵魂呼吸的通道。}
\verseline{}{阿耳肯斯之子站在华美的兵甲中，}
\verseline{}{披着针绣斗篷，伊比利亚的暗紫使他辉耀，}
\verseline{}{容貌出众；父亲阿耳肯斯把他送来，}
\verseline{}{他在叙迈托斯河畔的玛尔斯圣林中长大，}
\verseline{585}{那里有帕利科伊的丰厚祭坛，献祭足以安抚神祇。}
\verseline{}{墨泽提乌斯放下长枪，亲自把嘶鸣的投石索}
\verseline{}{以绷紧的皮带在头顶挥转三圈，}
\verseline{}{熔热的铅弹正面劈开少年的太阳穴，}
\verseline{}{把他伸展开来，铺在大片沙地上。\footnote{本卷第 569—589 行把杀人者、死者和兵器逐项连接：伊利俄纽斯以岩石杀琉刻提乌斯，利革尔杀厄玛提翁，阿西拉斯杀科律奈乌斯，凯纽斯杀俄耳提癸乌斯后又为图尔努斯所杀，卡皮斯的箭则在普里韦耳努斯伸手触伤时把手钉入左肋。墨泽提乌斯最后用投石索射出铅弹，击碎阿耳肯斯之子的太阳穴；诗句把高速铅弹写成仿佛熔热。}}
\verseline{590}{据说，阿斯卡尼乌斯这时第一次在战场上}
\verseline{}{拉弓射出疾箭；此前他只惯于惊吓奔逃的野兽，}
\verseline{}{如今却亲手击倒勇猛的努马努斯；}
\verseline{}{他又名雷穆卢斯；图尔努斯年纪较小的妹妹}
\verseline{}{新近才与他成婚，共处婚床。}
\verseline{595}{此人走在前列，向阵前高喊——有些值得转述，}
\verseline{}{有些不值一提；新得王亲的身份使他胸中骄矜，}
\verseline{}{以巨大的叫嚷夸示自己：}
\verseline{}{“难道不知羞耻？再一次被围困，困守壁垒，}
\verseline{}{两度陷落的佛里吉亚人，仍把城墙遮在死亡之前？}
\verseline{600}{看啊，这就是要在战争中求娶我们女子的人！}
\verseline{}{是哪位神把你们赶到意大利？又是哪一种疯狂？}
\verseline{}{这里没有阿特柔斯的儿子，也没有巧言造谎的乌利塞斯。}
\verseline{}{我们这一族生来强韧；婴儿一落地，}
\verseline{}{便被抱到河边，以刺骨寒冷和河水磨炼；}
\verseline{605}{少年彻夜打猎，使森林不得安宁，}
\verseline{}{他们以驯马为游戏，也以弯弓放箭为游戏。}
\verseline{}{青年能够忍受劳苦，惯于贫乏，}
\verseline{}{或以锄耙征服土地，或以战争震击城邑。}
\verseline{}{我们一生都消磨在钢铁中，又掉转枪杆，}
\verseline{610}{催逼年轻公牛的脊背；迟缓的老年}
\verseline{}{也不会削弱心灵力量，改变勇毅：}
\verseline{}{我们以头盔压住白发，永远乐于收集}
\verseline{}{新得的战利品，以抢夺所得为生。}
\verseline{}{你们却爱披藏红花与亮紫染成的彩衣，}
\verseline{615}{钟情懒散，乐于纵情舞蹈，}
\verseline{}{身穿有袖长袍，头巾又系着饰带。}
\verseline{}{啊，真正的佛里吉亚女人，绝非佛里吉亚男人！去吧，}
\verseline{}{去高耸的丁底马，那里双管笛吹出你们惯听的曲调！}
\verseline{}{伊达母神的手鼓和柏瑞钦托斯黄杨笛正在召唤你们。}
\verseline{620}{把兵器留给男人，向刀剑屈服吧。”\footnote{努马努斯又名雷穆卢斯，是新近迎娶图尔努斯妹妹的姻亲，和罗马起源传说中的罗慕路斯、雷穆斯没有关系。他以意大利人的河水磨炼、狩猎、农耕和战争生活贬斥佛里吉亚人的彩衣、舞蹈、长袖袍和母神乐器，又用“女人”否定对方的男子身份。}}
\verseline{}{他如此夸口，说出凶恶的话，}
\verseline{}{阿斯卡尼乌斯再也不能忍受；他迎面拉紧马鬃弦，}
\verseline{}{双臂向两侧分开，弯弓搭箭，}
\verseline{}{随后站定，先以誓愿恳求朱庇特：}
\verseline{625}{“全能的朱庇特，请对这大胆的事业点头应允。}
\verseline{}{我会亲自把依礼许下的祭品送到你的神庙，}
\verseline{}{在祭坛前安置一头额上镀金的小公牛；}
\verseline{}{它一身洁白，昂起的头已经与母牛等高，}
\verseline{}{已经能以牛角冲撞，以蹄子扬起沙土。”}
\verseline{630}{天父听见了，从晴朗天空的一侧}
\verseline{}{在左方发出雷鸣；致命的弓弦同时震响。}
\verseline{}{拉满的箭带着可怕的嘶声飞出，}
\verseline{}{直取雷穆卢斯的头颅，以铁锋贯穿}
\verseline{}{两边太阳穴。“去吧，以傲慢言辞嘲弄勇武！}
\verseline{635}{两度被攻陷的佛里吉亚人把这答复送给鲁图利人。”}
\verseline{}{阿斯卡尼乌斯只说这些。特乌克罗斯人以喊声相随，}
\verseline{}{在欢喜中咆哮，把心志高举到星辰。}
\verseline{}{那时，长发的阿波罗恰好坐在云上，}
\verseline{}{从高空的天域俯望奥索尼亚军阵与营城，}
\verseline{640}{他这样对获胜的尤卢斯说：}
\verseline{}{“少年，愿你初显的勇武日益增盛：人正由此走向星辰，}
\verseline{}{你由众神所生，也将生下众神。未来的一切战争}
\verseline{}{理当在阿萨剌科斯的族裔之下，依命运归于止息；}
\verseline{}{特洛伊已经容纳不下你。”他说完这些，}
\verseline{645}{便从高天投身而下，分开流动的风，}
\verseline{}{朝阿斯卡尼乌斯而去。随后，他把面容}
\verseline{}{变成老布忒斯；此人从前是达耳达诺斯人}
\verseline{}{安喀塞斯的执兵者，也是守卫门槛的忠仆，}
\verseline{}{后来父亲又令他伴随阿斯卡尼乌斯。阿波罗走来，}
\verseline{650}{一切都酷似这位老人：声音与肤色，}
\verseline{}{白发与声响森厉的兵甲；}
\verseline{}{他以这些话对斗志正烈的尤卢斯说：}
\verseline{}{“埃涅阿斯之子，你射杀努马努斯而未受报复，已经足够；}
\verseline{}{伟大的阿波罗准许你获得这第一次赞誉，}
\verseline{655}{并不嫉妒与你相同的弓箭；至于往后的战斗，}
\verseline{}{少年，暂且收手。”阿波罗这样开口时，}
\verseline{}{说到中途便离开凡人的视线，}
\verseline{}{远远从眼前消失，化入稀薄的空气。}
\verseline{}{达耳达诺斯首领认出了神明，也认出了神圣兵器，}
\verseline{660}{又从远去的声响中听出箭囊震鸣。}
\verseline{}{因此，他们依福玻斯的话与神威，制止了}
\verseline{}{渴望战斗的阿斯卡尼乌斯；自己却重新进入交锋，}
\verseline{}{把生命投向无遮无掩的危险。\footnote{阿斯卡尼乌斯第一次把惯用于狩猎的弓箭射向人；朱庇特以左侧雷鸣回应，阿波罗又赞许他的“新生勇武”和神裔未来。 “人正由此走向星辰”紧接少年第一次杀人，其含义首先属于战场上的功业与神裔未来。阿波罗随即化作安喀塞斯旧日执兵者布忒斯，命他以这第一次赞誉为足、退出其余战斗；赞许之后立即跟着制止。}}
\verseline{}{喊声沿整圈营墙的胸墙传开，}
\verseline{665}{战士拉紧劲弓，挥转投枪的皮带。}
\verseline{}{地面铺满投射兵器；盾牌和中空的头盔}
\verseline{}{被击得铿然作响，一场凶猛的战斗骤然升起：}
\verseline{}{如同在带雨的山羊星时节，从西方来的}
\verseline{}{暴雨以如此力量抽击大地；又如雷云}
\verseline{670}{把密集冰雹倾进海湾，正当朱庇特在南风中震怒，}
\verseline{}{卷动饱含水气的风暴，撕开天空中空的云层。}
\verseline{}{潘达鲁斯和比提亚斯是伊达山人阿尔卡诺耳之子，}
\verseline{}{林野女神雅埃拉在朱庇特的圣林里养大他们，}
\verseline{}{使青年身量比得上故乡山中的冷杉；\footnote{Conte 2019 在第 674 行采用 Bryant 的 in montibus；传本多作 et montibus。按此读法，abietibus ... patriis in montibus 意为“故乡山中的冷杉”，兄弟二人的身量与冷杉相比。}}
\verseline{675}{两人仗着一身武装，竟打开主帅交托守卫的营门，}
\verseline{}{主动邀请敌人进入壁垒。}
\verseline{}{他们自己站在门内，一左一右，守在塔前，}
\verseline{}{全身覆铁，高昂的头顶羽冠闪耀；}
\verseline{}{如同两株高入云端的橡树，清澈河水环流其旁，}
\verseline{680}{或立在波河岸边，或立在秀美的阿忒西斯河畔，}
\verseline{}{它们把未剪的树冠高高举向天空，}
\verseline{}{高耸的树梢随风摇动。\footnote{潘达鲁斯与比提亚斯是本卷守门兄弟，由林野女神雅埃拉养大；比提亚斯不同于第一卷中同名的迦太基贵族。波河与阿忒西斯河分别位于意大利北部，两株河岸橡树的明喻着重兄弟的身量、根基与高耸；开门决定的后果则由随后的战局展开。马克罗比乌斯（6.2.32）称，潘达鲁斯与比提亚斯开门迎战这一整段取自恩尼乌斯《编年史》第十五卷，其中两名希斯特里亚战士在围城中冲出城门、杀伤围军。维吉尔把早期罗马史诗的守城突击移入特洛伊人营地，并让开门最终反成图尔努斯闯入的灾祸。}}
\verseline{}{鲁图利人一看见入口敞开，便闯了进来；}
\verseline{}{奎耳肯斯、兵甲华美的阿奎库卢斯、}
\verseline{685}{性情急烈的特玛鲁斯和玛尔斯之裔海蒙}
\verseline{}{连同整个军阵，要么转身退却，}
\verseline{}{要么就在营门门槛上交出生命。}
\verseline{}{此时，敌对双方心中的怒火更盛；}
\verseline{}{特洛伊人也已经聚到同一处，结成大队，}
\verseline{690}{敢于近身交锋，又敢于向更远处冲出。}
\verseline{}{消息传到主帅图尔努斯那里；他正在另一边狂怒冲杀，}
\verseline{}{扰乱敌军，报信者告诉他：敌人}
\verseline{}{因新的屠杀沸腾，竟把大门敞开。}
\verseline{}{他抛下原来的行动，被可怕的怒火催动，}
\verseline{695}{猛扑向达耳达诺斯人的营门和傲慢的兄弟。}
\verseline{}{他先用投枪击倒安提法忒斯——}
\verseline{}{此人冲在最前，是高贵萨尔佩冬的私生子，}
\verseline{}{母亲来自忒拜。意大利山茱萸木飞过}
\verseline{}{轻柔的空气，刺进咽喉，深入胸膛；}
\verseline{700}{黑色伤洞喷出带沫的血流，铁锋}
\verseline{}{钉在肺中，渐渐被温热。}
\verseline{}{随后，他亲手击倒墨洛普斯和厄律玛斯，又杀死阿菲德努斯；}
\verseline{}{接着，他面对比提亚斯——双眼燃烧，心中咆哮——}
\verseline{}{却不用标枪，因为普通标枪不能取走他的生命；}
\verseline{705}{一支巨大的重型投枪旋转而来，发出尖啸，}
\verseline{}{如雷霆般被强力掷出；两层牛皮}
\verseline{}{和覆着双层鳞片与黄金的牢靠胸铠，都没能挡住它：}
\verseline{}{庞大的肢体一齐崩落，轰然坠下。}
\verseline{}{大地发出呻吟，巨盾在他身上震响。}
\verseline{710}{如同从前，拜亚那片优波亚人的海岸上，}
\verseline{}{一座石砌桥墩倒下；它以巨石先在陆上筑成，}
\verseline{}{再被投进海中，就这样向前倾覆，带着废墟}
\verseline{}{沉进深处，撞在海底；}
\verseline{}{海水相互翻搅，黑沙也高高涌起；}
\verseline{715}{高耸的普罗库塔因声响震颤，坚硬的床榻也震动；}
\verseline{}{那是伊纳里墨，奉朱庇特之命压在提丰身上。\footnote{图尔努斯杀比提亚斯所用 phalarica 是需要巨大力量投出的重型投枪；两层牛皮和金属鳞甲都未能抵挡。尸体倒下的明喻来自拜亚岸边筑港：巨石先在岸上砌成桥墩，再整体投入海中。接着又提到普罗库塔、伊纳里墨和被压在岛下的提丰，把坠落的震动写到坎帕尼亚一带。}}
\verseline{}{此时，执兵的玛尔斯把勇气与力量}
\verseline{}{添给拉丁人，在他们胸中转动尖锐的刺；}
\verseline{}{他又把“逃亡”与黑色的“恐惧”送入特乌克罗斯人中。}
\verseline{720}{敌人从各处聚来，因为交战的机会已经出现，}
\verseline{}{战神也降入他们心中。\footnote{玛尔斯在此亲自增强拉丁人的力量，又把“逃亡”与黑色的“恐惧”送入特乌克罗斯人中。“逃亡”与“恐惧”都被写成神明可以投下的力量；战神又直接“降入”士兵心中，使他们转而迎战。}}
\verseline{}{潘达鲁斯看见兄弟横尸在地，}
\verseline{}{明白命运已走到何处，何种变故正推动战局，}
\verseline{}{便以巨大力量转动门轴，关上营门，}
\verseline{725}{用宽阔的肩膀抵住它；许多自己人}
\verseline{}{仍被留在壁垒之外，在严酷的战斗中挣扎；}
\verseline{}{另一些人却被他同自己一并关入，迎进营内。}
\verseline{}{这个疯子没有看见鲁图利人的国王}
\verseline{}{正在队伍中央闯入，反倒亲手把他关进城中，}
\verseline{730}{如同把一头巨大的猛虎关入无力的牲畜之间。}
\verseline{}{立刻，一道异样的光从图尔努斯眼中迸出，}
\verseline{}{兵甲发出可怕的震响；血红的羽冠}
\verseline{}{在他头顶颤动，盾牌放射闪亮的雷光。}
\verseline{}{惊乱的埃涅阿斯族人认出了可恨的面容}
\verseline{735}{和庞大的四肢。高大的潘达鲁斯随即跃出，}
\verseline{}{因兄弟之死燃起怒火，开口说道：}
\verseline{}{“这里不是阿玛塔陪嫁的王宫，}
\verseline{}{也不是祖传城墙把图尔努斯围在其中的阿尔代亚。}
\verseline{}{你看见的是敌人的营地，再没有道路离开。”}
\verseline{740}{图尔努斯胸中平静，微笑着回答：}
\verseline{}{“若你心里还有勇武，就动手吧，与我交锋；}
\verseline{}{你会把消息带给普里阿摩斯：这里也找到了阿喀琉斯。”}
\verseline{}{他说完。潘达鲁斯使尽全力，旋臂掷出}
\verseline{}{一根满是木节、仍带粗糙树皮的长矛；}
\verseline{745}{风接住这道飞来的创伤；萨图尔努斯之女朱诺}
\verseline{}{飞来使它偏转，长矛钉进营门。\footnote{潘达鲁斯因兄弟倒下而关门，既把许多己方战士留在外面，也把图尔努斯关进营内。731行的异光可理解为从图尔努斯眼中闪出，也可联系众人眼前所显的光；兵甲、血红羽冠和盾面雷光一起使守军惊惶。潘达鲁斯以阿玛塔的嫁妆和阿尔代亚讥刺图尔努斯；图尔努斯则自比阿喀琉斯，声称潘达鲁斯将向普里阿摩斯报信。朱诺使潘达鲁斯掷出的长矛偏转，钉入营门。}}
\verseline{}{“可你逃不过我右手全力挥出的这一击；}
\verseline{}{挥动兵器的人不同，造成的伤也不同。”}
\verseline{}{说着，图尔努斯高高举剑，整个人随之挺起，}
\verseline{750}{以铁刃从两边太阳穴之间劈开额头，}
\verseline{}{劈开尚未生须的双颊，留下巨大的伤口。}
\verseline{}{巨声轰响，大地被庞大重量震动；}
\verseline{}{潘达鲁斯临死，把倒下的四肢和沾满脑浆的兵甲}
\verseline{}{铺在地上；他的头颅正中劈成两半，}
\verseline{755}{从两侧肩头一边一半垂落。}
\verseline{}{特洛伊人背转身体，在惊惶恐惧中四散；}
\verseline{}{倘若胜利者立即想到}
\verseline{}{亲手打破门闩，让同伴从大门涌入，}
\verseline{}{那一天本可成为战争与特洛伊民族的最后一天；}
\verseline{760}{可是狂暴与疯狂的杀戮欲驱使着怒火正炽的他，}
\verseline{}{直冲那些已经转身逃跑的人。\footnote{叙述者明确给出反事实判断：图尔努斯若在杀死潘达鲁斯后立即开门放入同伴，那一天可能成为战争与特洛伊民族的末日。他却受 furor 和眼前的杀戮欲驱使，没有开门，随后独自追杀逃兵。Conte 2019 正文采用 Bentley 的 auersos（“转身背向”），传本多作 aduersos（“迎面相对”）；前者与前文“特洛伊人转身逃散”以及随后“从背后”掷枪相互呼应。}}
\verseline{}{他先杀死法勒里斯，又割断古革斯的膝弯；}
\verseline{}{随后夺过投枪，掷向奔逃者}
\verseline{}{的背部；朱诺供给他力量和心志。}
\verseline{765}{他又杀了哈吕斯，也杀死盾牌已被钉穿的斐革乌斯，}
\verseline{}{接着杀死墙上那些全不知情、仍在唤起战事的人：}
\verseline{}{阿尔坎得尔、哈利乌斯、诺厄蒙和普律塔尼斯。}
\verseline{}{林叩斯正迎面赶来，大声召唤同伴，}
\verseline{}{图尔努斯从土垒右侧俯身，以闪亮的剑}
\verseline{770}{抢先击中；只消近身一击，}
\verseline{}{林叩斯的头连同头盔被斩下，落在远处。随后，}
\verseline{}{图尔努斯又杀了屠猎野兽的阿米库斯；无人比他更善于}
\verseline{}{亲手涂抹投枪，让铁锋带上毒液。}
\verseline{}{又杀死埃俄卢斯之裔克吕提乌斯，以及缪斯所爱的克瑞透斯；}
\verseline{775}{克瑞透斯是缪斯的同伴，永远钟情诗歌}
\verseline{}{与竖琴，钟情于让节律在琴弦上绷紧。}
\verseline{}{他总在歌唱战马、武器、英雄与战斗。\footnote{本卷第 762—777 行继续列出图尔努斯所杀之人，又在克瑞透斯身上突然停留：他是缪斯的朋友，爱诗、琴弦与节律，也总歌唱战马、兵器、英雄和战斗。诗人把一位战争歌者写进战斗名录。}}
\verseline{}{特乌克罗斯人的统帅终于听说己方遭到屠杀，}
\verseline{}{门涅斯透斯和勇烈的塞雷斯图斯一同赶来，}
\verseline{780}{看见同伴四散，敌人已被接入营内。}
\verseline{}{门涅斯透斯说：“你们还要逃向何处？还要奔往哪里？}
\verseline{}{你们还有什么城墙，前面还有哪一道壁垒？}
\verseline{}{同胞们，一个人被你们自己的土垒}
\verseline{}{从四面包围，竟能在营城内安然造成如此多的屠杀，}
\verseline{785}{把这么多最优秀的青年送入奥尔库斯？}
\verseline{}{懦弱的人啊，你们既不怜悯不幸的祖国和古老的神明，}
\verseline{}{也不因伟大的埃涅阿斯感到羞耻吗？”}
\verseline{}{这些话点燃众人；他们坚定下来，以密集队形}
\verseline{}{站稳。图尔努斯渐渐退出战斗，}
\verseline{790}{走向河流和被水环绕的一侧。}
\verseline{}{特乌克罗斯人见状，更猛烈地压上，发出巨大的喊声，}
\verseline{}{把队伍聚成一团。如同一群猎人以敌对的兵器}
\verseline{}{紧逼一头凶猛的狮子；它虽然受惊，}
\verseline{}{仍粗暴地怒视，向后退去；怒气与勇武}
\verseline{795}{不容它露出脊背，可它虽想向前，}
\verseline{}{却无法穿过密集的兵器和战士：}
\verseline{}{迟疑的图尔努斯也正这样，不疾不徐地}
\verseline{}{倒退脚步，心中因愤怒翻腾。}
\verseline{}{即使如此，他当时仍两次冲入敌阵中央，}
\verseline{800}{两次迫使混乱的军阵沿墙奔逃；}
\verseline{}{但整个营地的军力迅速聚合为一，}
\verseline{}{萨图尔努斯之女朱诺也不敢再为他供给力量，}
\verseline{}{因为朱庇特已从天空派下伊里斯，}
\verseline{}{给他的姊妹带来毫不温和的命令：}
\verseline{805}{图尔努斯必须退出特乌克罗斯人的高墙。}
\verseline{}{因此，这青年无论用盾牌}
\verseline{}{还是右手，都再也无力承受如此重压；四面投来的兵器}
\verseline{}{正把他压倒。头盔在太阳穴四周不断铿鸣，}
\verseline{}{坚固的青铜在石块下裂开，}
\verseline{810}{头上的羽饰也被击落；盾心承受不住}
\verseline{}{一次次打击；特洛伊人不断投枪，门涅斯透斯}
\verseline{}{也如闪电般加倍猛攻。这时，汗水从他全身}
\verseline{}{淌落，汇成沥青般黏黑的流瀑；他再不能喘息，}
\verseline{}{痛苦的喘息使疲惫的四肢阵阵震颤。}
\verseline{815}{直到这时，他才带着全副武装，纵身}
\verseline{}{从高处投入河中；河流以浑黄的漩涡}
\verseline{}{接住落下的他，以柔软的波浪托起，}
\verseline{}{洗去屠杀的血污，把欣喜的他送还给同伴。\footnote{门涅斯透斯重新集结守军后，图尔努斯仍两次冲散敌阵；直到朱庇特派伊里斯向朱诺传下严厉命令，朱诺才不敢继续补足他的力量。他在盾、盔、石块、枪矛、汗水和喘息的共同压力下跳入台伯河。河水以柔波接住他，洗去他身上的屠杀血污，把获救而欣喜的图尔努斯送回同伴，战争仍将继续。马克罗比乌斯（6.3.2—4）指出，《伊利亚特》先写大埃阿斯受密集投枪压迫，恩尼乌斯《编年史》第十五卷又以相近的战斗描写写罗马军官盖乌斯·埃利乌斯，维吉尔这里再用于被困营中的图尔努斯；盔鸣、盾击、汗流和不能喘息等细节都可比较。}}
